%%%%%%%%%%%%%%%%%%%%%%%%%%%%%%%%%%%%%%%%%%%%%%%%%%%%%%%%%%%%%%%%%
                % Probabilistic transformer %
%%%%%%%%%%%%%%%%%%%%%%%%%%%%%%%%%%%%%%%%%%%%%%%%%%%%%%%%%%%%%%%%%


As mentioned in chapter 1, there is a theoretical gap in the area. {\color{orange} está bien dicho theoretical gap?}
A study of non-determinism in transformers is missing.

When $\TF$ is a deterministic transformer, $\TF^i(\alpha)$ is an exact symbol for all $\alpha \in \Sigma^*$ and $i \in \mathbb{N}$.
In counterpart, when the token selection process of the transformer samples the distribution, it becomes a random variable.
Therefore, we introduce the following definition.


\begin{definition}[Probabilistic transformer]
    We call probabilistic transformers to the transformers that sample the distribution for selecting the next token.
    We say that a family of probabilistic transformers $\{\TF_n\}_{n \in \mathbb{N}}$ recognize a language $\gL$ when: 
    \[
    \alpha \in \gL \iff \Pr[\TF_{|\alpha|}^*(\alpha) = \qaccept] > 2/3
    \]
\end{definition}


